\section*{Capítulo \ref{ch:g10:relative-motion} --- Movimento relativo}

\begin{solution}{exo:g10:relative-motion:1}
Em relação ao vagão: em repouso. Em relação ao solo e à árvore: em
movimento a \qty{100}{km/h}. Em relação ao trem que vem em sentido
contrário: velocidades em sentidos opostos se somam,
$100 + 100 = \qty{200}{km/h}$. No \omterm{def:g10:relative-motion:frame}{referencial} do trem do passageiro, a
árvore se move para trás a \qty{100}{km/h} ao longo de uma reta.
\end{solution}

\begin{solution}{exo:g10:relative-motion:2}
$72 / 3.6 = \qty{20}{m/s}$; $108 / 3.6 = \qty{30}{m/s}$;
$25 \times 3.6 = \qty{90}{km/h}$; $340 \times 3.6 = \qty{1224}{km/h}$. Em
cada par o valor em \unit{km/h} é o maior, como a regra de conferência
exige.
\end{solution}

\begin{solution}{exo:g10:relative-motion:3}
Ciclista: $v = \frac{\qty{36}{km}}{\qty{1.5}{h}} = \qty{24}{km/h}
= 24/3.6 \approx \qty{6.7}{m/s}$. Velocista:
$v = \frac{\qty{100}{m}}{\qty{10.0}{s}} = \qty{10.0}{m/s}
= \qty{36}{km/h}$. O velocista é uma vez e meia mais rápido --- mas só
por dez segundos, enquanto o ciclista mantém o ritmo por uma hora e meia.
\omterm{def:g10:relative-motion:speed}{Velocidades médias} só se comparam com justiça em durações comparáveis.
\end{solution}

\begin{solution}{exo:g10:relative-motion:4}
Mesmo sentido: $1.4 + 0.9 = \qty{2.3}{m/s}$. Contra a esteira:
$1.4 - 0.9 = \qty{0.5}{m/s}$. Parado sobre a esteira: os \qty{0.9}{m/s}
da própria esteira.
\end{solution}

\begin{solution}{exo:g10:relative-motion:5}
Caneta solta: o \omterm{def:g10:relative-motion:frames}{referencial do solo} funciona, mas o \omterm{def:g10:relative-motion:frame}{referencial} do
\emph{trem} é ainda melhor --- nele a caneta simplesmente cai na
vertical. Estação espacial: \omterm{def:g10:relative-motion:frames}{referencial geocêntrico} (ela orbita a Terra
como um todo, ignorando a rotação). Ano marciano: \omterm{def:g10:relative-motion:frames}{referencial
heliocêntrico} (Marte descreve um quase-círculo em volta do Sol). Corrida:
\omterm{def:g10:relative-motion:frames}{referencial do solo} (largada, chegada e cronômetro estão todos fixos ao
solo).
\end{solution}

\begin{solution}{exo:g10:relative-motion:6}
\emph{1.} $108/3.6 = \qty{30}{m/s}$. Para a frente:
$30 + 1.5 = \qty{31.5}{m/s}$; para trás:
$30 - 1.5 = \qty{28.5}{m/s}$.

\emph{2.} Em relação ao trem, o passageiro percorre \qty{60}{m} a
\qty{1.5}{m/s}: $t = 60/1.5 = \qty{40}{s}$. Em relação ao solo, ele
percorre $31.5 \times 40 = \qty{1260}{m}$ --- os \qty{1200}{m}
transportados pelo trem mais os \qty{60}{m} caminhados, como a regra de
composição prevê.
\end{solution}

\begin{solution}{exo:g10:relative-motion:7}
\emph{1.} Rio abaixo: $5.0 + 2.0 = \qty{7.0}{m/s}$, logo
$t_1 = 420/7.0 = \qty{60}{s}$. Rio acima: $5.0 - 2.0 = \qty{3.0}{m/s}$,
logo $t_2 = 420/3.0 = \qty{140}{s}$.

\emph{2.} \omterm{def:g10:relative-motion:speed}{Velocidade média}: $v = \frac{840}{60 + 140}
= \frac{840}{200} = \qty{4.2}{m/s} < \qty{5.0}{m/s}$. A ajuda e o
estorvo da correnteza são iguais em \emph{velocidade}, mas não em
\emph{tempo}: o barco passa \qty{140}{s} sendo freado e apenas
\qty{60}{s} sendo ajudado, de modo que o trecho lento domina a média.
\end{solution}

\begin{solution}{exo:g10:relative-motion:8}
\emph{1.} Sentidos opostos: no \omterm{def:g10:relative-motion:frame}{referencial} de A, B se aproxima a
$25 + 20 = \qty{45}{m/s}$. O cruzamento se completa quando B tiver
deslizado pelo comprimento combinado $180 + 90 = \qty{270}{m}$:
$t = 270/45 = \qty{6.0}{s}$.

\emph{2.} Mesmo sentido: velocidade relativa $25 - 20 = \qty{5.0}{m/s}$,
e $t = 270/5.0 = \qty{54}{s}$ --- nove vezes mais, e é por isso que
ultrapassar um caminhão parece não acabar nunca.
\end{solution}

\begin{solution}{exo:g10:relative-motion:9}
\omterm{def:g10:relative-motion:frame}{Referencial} do ciclista: um círculo (o círculo do aro, centrado no cubo).
\omterm{def:g10:relative-motion:frames}{Referencial do solo}: uma \omterm{prop:g10:relative-motion:relativity}{cicloide} --- uma sequência de arcadas, uma por
volta da roda. \omterm{def:g10:relative-motion:frame}{Referencial} da própria válvula: a válvula está em repouso,
um único ponto. O cubo no \omterm{def:g10:relative-motion:frames}{referencial do solo}: uma reta horizontal à
altura de um raio de roda.
\end{solution}

\begin{solution}{exo:g10:relative-motion:10}
\emph{1.} $v = \frac{\qty{40000}{km}}{\qty{24}{h}}
\approx \qty{1.7e3}{km/h} = 1667/3.6 \approx \qty{463}{m/s}$ --- mais
rápido do que o som.

\emph{2.} Não a sentimos porque tudo à nossa volta --- solo, ar,
edifícios --- compartilha exatamente o mesmo movimento, de modo que nada
se move em relação a nós (e o princípio da inércia,
\cref{ch:g10:inertia}, explicará por que um movimento uniforme
compartilhado é indetectável por dentro). ``O Sol nasce'' é o relato no
\omterm{def:g10:relative-motion:frames}{referencial do solo} e ``a Terra gira'' é o relato no \omterm{def:g10:relative-motion:frames}{referencial
geocêntrico} de uma única e mesma observação.
\end{solution}

\begin{solution}{exo:g10:relative-motion:11}
Seja $L$ o comprimento da escada. Velocidade da escada: $L/60$;
velocidade de caminhada: $L/90$. Caminhando na escada em movimento, as
velocidades se somam:
\[
v = \frac{L}{60} + \frac{L}{90} = L\left(\frac{3 + 2}{180}\right)
= \frac{L}{36},
\]
de modo que a subida leva \qty{36}{s}. (Não é a média de $60$ e $90$: as
velocidades se somam, os tempos não.)
\end{solution}

\begin{solution}{exo:g10:relative-motion:12}
\emph{1.} Para uma distância $d$ em cada sentido:
$t = \frac{d}{60} + \frac{d}{90} = d\,\frac{3 + 2}{180} = \frac{d}{36}$,
logo $v = \frac{2d}{d/36} = \qty{72}{km/h}$.

\emph{2.} O ingênuo $75$ tira a média das duas velocidades como se o
motorista tivesse passado \emph{tempos} iguais em cada uma; na verdade as
distâncias é que são iguais, de modo que o trecho mais lento dura mais
($\frac{d}{60} > \frac{d}{90}$) e puxa a média para perto de
\qty{60}{km/h}. A \omterm{def:g10:relative-motion:speed}{velocidade média} é a distância total dividida pelo
tempo total --- e nada mais.
\end{solution}

\begin{solution}{exo:g10:relative-motion:13}
\emph{1.} No \omterm{def:g10:relative-motion:frame}{referencial} do primeiro ciclista, o segundo se aproxima em
linha reta a $15 + 20 = \qty{35}{km/h}$, a partir de \qty{35}{km} de
distância: eles se encontram depois de exatamente \qty{1.0}{h}.

\emph{2.} O primeiro pedalou $15 \times 1 = \qty{15}{km}$ e o segundo,
$20 \times 1 = \qty{20}{km}$ --- ao todo \qty{35}{km}, como tinha de ser.

\emph{3.} Mesmo sentido: velocidade de aproximação
$20 - 15 = \qty{5}{km/h}$, logo a perseguição dura
$35/5 = \qty{7}{h}$.
\end{solution}

\begin{solution}{exo:g10:relative-motion:14}
\emph{1.} A correnteza empurra o nadador ao longo do rio, não através
dele: os dois movimentos se compõem de modo independente, e o avanço em
direção à outra margem se faz a \qty{1.2}{m/s} qualquer que seja a
deriva. Tempo de travessia: $t = 60/1.2 = \qty{50}{s}$.

\emph{2.} Deriva: $0.9 \times 50 = \qty{45}{m}$ rio abaixo.

\emph{3.} A velocidade no \omterm{def:g10:relative-motion:frames}{referencial do solo} é a hipotenusa de um
triângulo retângulo de catetos $1.2$ e $0.9$:
$v = \sqrt{1.2^2 + 0.9^2} = \sqrt{2.25} = \qty{1.5}{m/s}$ (um triângulo
$3$--$4$--$5$ multiplicado por $0.3$). Comprimento do caminho:
$1.5 \times 50 = \qty{75}{m}$ --- ou diretamente
$\sqrt{60^2 + 45^2} = \qty{75}{m}$.
\end{solution}

\begin{solution}{exo:g10:relative-motion:15}
\emph{1.} $v = \frac{2\pi \times \num{1.50e11}}{\num{3.16e7}}
= \frac{\num{9.42e11}}{\num{3.16e7}} \approx \qty{2.98e4}{m/s}
\approx \qty{30}{km/s}$.

\emph{2.} $30 / 0.46 \approx 65$: a Terra nos leva em volta do Sol cerca
de sessenta e cinco vezes mais depressa do que nos gira em torno do seu
eixo.

\emph{3.} No \omterm{def:g10:relative-motion:frames}{referencial geocêntrico}, o caminho aparente de Marte combina
o movimento orbital dele com o da Terra e, quando a Terra, mais rápida,
ultrapassa Marte, a combinação corre para trás por algumas semanas. No
\omterm{def:g10:relative-motion:frames}{referencial heliocêntrico} os dois planetas simplesmente circulam o Sol em
ritmos regulares --- a ``volta para trás'' não é um movimento de Marte,
mas um artefato do \omterm{def:g10:relative-motion:frame}{referencial}.
\end{solution}

\begin{solution}{pb:g10:relative-motion:1}
\textbf{1.} $t = 90/1.0 = \qty{90}{s}$.

\textbf{2.} Velocidade em relação ao solo $1.5 + 1.0 = \qty{2.5}{m/s}$:
$t = 90/2.5 = \qty{36}{s}$.

\textbf{3.} $t = 90/1.5 = \qty{60}{s}$.

\textbf{4.} $1.5 - 1.0 = \qty{0.5}{m/s}$ em relação ao solo, contra o
sentido da esteira: $t = 90/0.5 = \qty{180}{s}$ --- três minutos
inteiros de caminhada acelerada.

\textbf{5.} Velocidade em relação ao solo $1.0 - 1.0 = \qty{0}{m/s}$: a
criança corre com afinco e não sai do lugar, em repouso no \omterm{def:g10:relative-motion:frames}{referencial do
solo} enquanto dispara no \omterm{def:g10:relative-motion:frame}{referencial} da esteira. Isso é uma esteira
ergométrica.

\textbf{6.} \omterm{def:g10:relative-motion:frame}{Referencial} da esteira / \omterm{def:g10:relative-motion:frames}{referencial do solo}: viajante parado
$0$ / \qty{1.0}{m/s}; caminhante a favor da esteira $1.5$ /
\qty{2.5}{m/s}; caminhante no piso $0.5$ (o piso desliza para trás no
\omterm{def:g10:relative-motion:frame}{referencial} da esteira) / \qty{1.5}{m/s}; brincalhão $1.5$ /
\qty{0.5}{m/s}; criança $1.0$ / \qty{0}{m/s}. As pernas de cada um sentem
a coluna do \omterm{def:g10:relative-motion:frame}{referencial} daquilo que está sob os seus pés --- a coluna da
esteira para quem anda sobre ela, a do solo para quem anda no piso: o
esforço é a velocidade de caminhada em relação ao que está por baixo.

\textbf{7.} O avanço através do rio se faz a \qty{2.0}{m/s} faça a
correnteza o que fizer: $t = 120/2.0 = \qty{60}{s}$.

\textbf{8.} Deriva: $1.5 \times 60 = \qty{90}{m}$ rio abaixo.

\textbf{9.} Velocidade em relação ao solo: hipotenusa de um triângulo
retângulo de catetos $2.0$ e $1.5$: $v = \sqrt{2.0^2 + 1.5^2}
= \sqrt{6.25} = \qty{2.5}{m/s}$. Comprimento da \omterm{def:g10:relative-motion:trajectory}{trajetória}:
$2.5 \times 60 = \qty{150}{m}$ --- o trio $3$--$4$--$5$ de novo,
ampliado para $90$--$120$--$150$.

\textbf{10.} \omterm{def:g10:relative-motion:frame}{Referencial} da água: uma reta de \qty{120}{m},
perpendicular às margens. \omterm{def:g10:relative-motion:frames}{Referencial do solo}: também uma reta, mas
inclinada e com \qty{150}{m}. Mesma forma (reta), direção e comprimento
diferentes: a \omterm{def:g10:relative-motion:trajectory}{trajetória} pertence ao \omterm{def:g10:relative-motion:frame}{referencial}, e não à balsa.

\textbf{11.} O avanço através do rio se faz a \qty{1.0}{m/s}, logo
$t = 120/1.0 = \qty{120}{s}$. Deriva:
$1.5 \times 120 = \qty{180}{m}$ --- o dobro dos \qty{90}{m} da balsa.
Velocidade em relação ao solo: hipotenusa de catetos $1.0$ e $1.5$:
$v = \sqrt{1.0^2 + 1.5^2} = \sqrt{3.25} \approx \qty{1.8}{m/s}$. O
nadador, mais lento, demora o dobro para atravessar, de modo que a mesma
correnteza tem o dobro de tempo para levá-lo rio abaixo.

\textbf{12.} \omterm{def:g10:relative-motion:frames}{Referencial do solo}: o tronco deriva com os \qty{1.5}{m/s}
da correnteza. \omterm{def:g10:relative-motion:frame}{Referencial} da água: ele está em repouso, boiando com a
água à sua volta. Se o tronco ``se move'' não tem resposta independente
de \omterm{def:g10:relative-motion:frame}{referencial} --- e é esse o ponto do capítulo inteiro.

\textbf{13.} O \omterm{def:g10:relative-motion:frames}{referencial geocêntrico}. Raio da \omterm{def:g10:universal-gravitation:satellite}{órbita}
$r = \num{6.37e6} + \num{4.0e5} = \qty{6.77e6}{m}$, logo
\[
v = \frac{2\pi r}{T} = \frac{2\pi \times \num{6.77e6}}{\num{5.56e3}}
\approx \qty{7.65e3}{m/s} \approx \qty{7.7}{km/s}.
\]

\textbf{14.} Enquanto a estação completa uma \omterm{def:g10:universal-gravitation:satellite}{órbita}, a Terra gira por
baixo dela cerca de $23^\circ$, de modo que o rastro sobre o solo se
desloca para oeste a cada passagem. No \omterm{def:g10:relative-motion:frames}{referencial do solo} o caminho da
estação é uma curva ondulada complicada que nunca se repete exatamente;
no \omterm{def:g10:relative-motion:frames}{referencial geocêntrico} é um único quase-círculo fixo.

\textbf{15.} $7.65 / 0.46 \approx 17$: a estação supera o equador em
rotação em cerca de dezessete vezes, e é por isso que ela cruza o céu
inteiro em minutos.

\textbf{16.} $v = \frac{2\pi \times \num{4.22e7}}{\num{8.64e4}}
\approx \qty{3.07e3}{m/s} \approx \qty{3.1}{km/s}$ no \omterm{def:g10:relative-motion:frames}{referencial
geocêntrico}. No \omterm{def:g10:relative-motion:frames}{referencial do solo} ele está \emph{em repouso}, pairando
sobre um ponto fixo do equador --- ele dá a volta no eixo da Terra
exatamente no mesmo tempo que o próprio solo, de modo que os dois
movimentos se cancelam no \omterm{def:g10:relative-motion:frames}{referencial do solo}. Em repouso e a
\qty{3.1}{km/s}: as duas coisas verdadeiras, uma por \omterm{def:g10:relative-motion:frame}{referencial}.

\textbf{17.} No \omterm{def:g10:relative-motion:frames}{referencial heliocêntrico} a Lua percorre essencialmente a
\omterm{def:g10:universal-gravitation:satellite}{órbita} da Terra --- um enorme quase-círculo em volta do Sol ---
sobreposta a uma oscilação suave de ciclo de \qty{27}{d}. A influência do
Sol domina de modo tão forte que o caminho se curva sempre em direção ao
Sol; ele nunca dá voltas para trás.

\textbf{18.} ``O Sol nasce'' é o relato no \omterm{def:g10:relative-motion:frames}{referencial do solo}, onde o
observador está em repouso e o Sol varre o céu. ``A Terra gira'' é o
relato no \omterm{def:g10:relative-motion:frames}{referencial geocêntrico}, onde a direção do Sol é fixa e o
observador é levado para leste em volta do eixo. As duas descrevem a
mesma observação; cada uma está correta no seu \omterm{def:g10:relative-motion:frame}{referencial}, e a
\cref{prop:g10:relative-motion:relativity} diz que nenhuma observação do
movimento sozinha pode arbitrar.

\textbf{19.} O \omterm{def:g10:relative-motion:frame}{referencial} de Ptolomeu, centrado na Terra, casava
naturalmente com o movimento diário do céu e previa as posições
planetárias de forma utilizável, mas só ao preço de voltas sobre voltas
(epiciclos) para reproduzir as andanças de cada planeta. O \omterm{def:g10:relative-motion:frame}{referencial} de
Copérnico, centrado no Sol, fez de cada planeta um quase-círculo regular
e transformou as voltas para trás de Marte num efeito de \omterm{def:g10:relative-motion:frame}{referencial} ---
a Terra ultrapassando um corredor mais lento pela raia de dentro. O
veredito moderno não é que o Sol ``realmente'' fique parado (ele também
se move, em volta do centro da galáxia), mas que referenciais são
ferramentas: o heliocêntrico torna o movimento planetário
incomparavelmente mais simples, e é para a simplicidade, não para a
verdade, que serve um \omterm{def:g10:relative-motion:frame}{referencial}.

\textbf{20.} \omterm{def:g10:relative-motion:frames}{Referencial do solo}: \qty{0}{km/s} --- em repouso na
poltrona. \omterm{def:g10:relative-motion:frames}{Referencial geocêntrico}: cerca de \qty{0.46}{km/s}, levado em
volta do eixo da Terra. \omterm{def:g10:relative-motion:frames}{Referencial heliocêntrico}: cerca de
\qty{30}{km/s}, levado em volta do Sol. Durante um filme de
\qty{90}{min}: $d = \num{3.0e4} \times 5400 \approx \qty{1.6e8}{m}
= \qty{160000}{km}$ --- quase metade do caminho até a Lua sem sair da
poltrona. Moral: o movimento não é uma propriedade de um corpo, mas de um
corpo e de um \omterm{def:g10:relative-motion:frame}{referencial}; escolha o \omterm{def:g10:relative-motion:frame}{referencial} que torna a história
simples e diga sempre qual foi.
\end{solution}
